\documentclass{article}
\usepackage{amssymb}
\usepackage{courier}
\usepackage{fancyhdr}
\usepackage{fancyvrb}
\usepackage[top=.75in, bottom=.75in, left=.75in,right=.75in]{geometry}
\usepackage{graphicx}
\usepackage{lastpage}
\usepackage{listings}
\lstset{basicstyle=\small\ttfamily}
\usepackage{soul}
\usepackage{tabularx}
\usepackage{upquote}
\usepackage{xcolor}

\usepackage[colorlinks,urlcolor={blue}]{hyperref}

\begin{document}

\fancyfoot[L]{\color{gray} C4CS -- W'16}
\fancyfoot[R]{\color{gray} Revision 1.0}
\fancyfoot[C]{\color{gray} \thepage~/~\pageref*{LastPage}}
\pagestyle{fancyplain}



\title{\textbf{Homework 4\\Git}}
\author{Assigned: Friday, January 29, 11:00AM}
\date{\textbf{\color{red}{Due: Friday, February 5, 11:00AM (Hard Deadline)}}}
\maketitle


\section*{Submission Instructions}
Submit this assignment on \href{https://gradescope.com/courses/2248}{Gradescope}.
You must submit every page of this PDF.
We recommend using the free online tool \href{https://www.pdfescape.com}{PDFescape}
to edit and fill out this PDF.
You may also print, handwrite, and scan this assignment.

\medskip
\noindent
\hl{Even if there are not questions to answer on the first page, please still
submit every page of this PDF.}

\medskip
\noindent
There may multiple answers for each question. If you are unsure,
state your assumptions and your reasoning for why you think your answer
makes sense.

\section{Git's Chicken and Egg Problem}

Git has a (well-deserved) reputation as a tool with a high learning curve.
There are several reasons for this. One is that git has a bad user interface
-- the commands and flags you use to get things done are frustratingly
inconsistent and simply must be learned over time.
The bigger reason, however, is that using git requires having an understanding
of how git works and learning how git works requires having some experience
using git. This circular dependency makes it hard to get started as there is
no good way to break in. The best we can do is try to hold your hand as you
get started.

\medskip
\noindent
Code Academy and GitHub teamed up to build a nice, online git beginners tool.
As you work through the exercises, pay attention to the explanations of each
step at the top of the page, and the ``Advice'' section at the
bottom/bottom-right of the page. The text introduces several new pieces of
git-specific language.

\medskip
\noindent
\textbf{Complete the online course at \url{https://try.github.io}. Grab a
  screenshot of your badge and copy the URL here.
  {\small (You will need to sign in or create an account get your badge)}
}

\medskip
\noindent
\fbox{
  \begin{minipage}{0.5\textwidth}
    My badge and URL as an example:

    \includegraphics[width=\textwidth]{badge}

    \footnotesize
    \url{https://www.codeschool.com/users/2276474/badges/121}
  \end{minipage}
  \begin{minipage}{0.5\textwidth}
    \begin{center}
      Your badge here
    \end{center}
  \end{minipage}
}


\subsection*{Some other tutorials / resources}

There is a \emph{lot} written on the Internet about how to learn git. Some
highlights:

\begin{itemize}
  \item For a deeper dive into git from the basics to the very advanced, check
    out \href{https://www.atlassian.com/git/tutorials}{Atlassian's tutorials}
    (Atlassian is the company behind BitBucket and JIRA).
  \item For more visual learners, check out this
    \href{http://pcottle.github.io/learnGitBranching/}{interactive visual
    learning game} -- caveat this game simplifies some things, but I think
    it's still a useful tool to learn from.  There's even some
    \href{https://github.com/pcottle/learnGitBranching/issues/201}{low-hanging
    fruit} lying around to improve this tool.
\end{itemize}


\newpage
\section{Using git in your own projects}
\textbf{\color{red} This question sets up the Winter \ul{2015} EECS 280 P2 as
an example. Don't use it for this term!!}
\begin{enumerate}
    \setcounter{enumi}{-1}
  \item Install git: \texttt{sudo apt-get install git}
\end{enumerate}
\noindent
Using version control well is a habit learned over time. As a first step, for
the rest of the term, the \textbf{first} thing you should do whenever an EECS
class assigns you a project is:
\begin{enumerate}
  \item Create a repository
    \begin{Verbatim}[fontsize=\footnotesize]
    > mkdir -p eecs280-w15/project2
    > cd $_
    > git init
    \end{Verbatim}
  \item Grab a copy of the spec
    \begin{Verbatim}[fontsize=\footnotesize]
    > wget 'https://drive.google.com/uc?id=0B4qlH840ZwikRmQtdkRIeTZjODA&export=download' -O spec.pdf
    \end{Verbatim}
  \item Grab any starter code
    \begin{Verbatim}[fontsize=\footnotesize]
    > wget 'https://drive.google.com/uc?id=0B4qlH840ZwikbkZLS3Z5YTVSeW8&export=download' -O eecs280-w15-p2.tgz
    > tar -xf eecs280-w15-p2.tgz
    > rm $_
    \end{Verbatim}
  \item Add it all to git and commit it
    \begin{lstlisting}
    > git add *
    > git status # Check to make sure everything looks right
    > git commit
    \end{lstlisting}
\end{enumerate}
Do this \textbf{before} even beginning to read the spec. Make it a habit.

\medskip
\noindent
If you try typing \texttt{make}, you'll find that this starter code does not
compile. The spec has a great tip at the bottom of page~11 that suggests
adding stub functions.

\medskip
\noindent
There are 15 functions we need to stub out. That could be a pain, but
fortunately the header file has already done most of the work. We'll drill
into this command next week, but for today, simply run:
\begin{lstlisting}
> echo -e '#include <cassert>\n#include "p2.h"\n' > p2.cpp
> grep ';' p2.h | grep -v ' \*' >> p2.cpp
\end{lstlisting}

\medskip
\noindent
One of the optional things to learn about editors from last week's homework
was macros. In case you missed it, let's demonstrate their use here:\\
{\small \emph{
    Note: I use vim, but emacs also supports macros with \texttt{C-x (},
    \texttt{C-x )}, and \texttt{C-x e} to record, finish, and execute.
}}
\begin{lstlisting}
> vim p2.cpp
# press 'jjj'    so that your cursor is on the i of int sum(...
# press 'qq'     to begin recording a macro into the vim register q
# press 'f;'     jump to the ';' character
# press 's'      remove character under the cursor and enter insert mode
# insert the needed text:
#   '<space>{<enter>assert(false);<enter>}<enter><escape>'
# press 'j'      so that your cursor is on the i of int product(...
# press 'q'      to finish recording the macro
# press '16@q'   14 times play the commands stored in register q
# save and quit vim
\end{lstlisting}
\noindent

\noindent
Now try \texttt{make}. Cool, a building project in 60\,seconds or less.

\medskip
\noindent
\textbf{Commit the changes.}

\medskip
\noindent
You ran \texttt{git status} before committing those changes, right? (Because
it's \emph{always} a good idea to run \texttt{git status} \emph{all} the time).
You probably noticed that annoying ``Untracked Files'' section that
lists all of the compiled output the Makefile made. It doesn't really make
sense to track compiled output in a version control system (it is hard to
track changes on a binary file such as an executable).\footnote{
  Some binary files are good to check in though, such as the project
  specification, which will rarely or never change.
}
Instead, we would prefer if git would \hl{ignore} the built output.

\medskip
\noindent
\textbf{Commit a change so that built output is ignored.} When you are done,
\texttt{git status} should print
\begin{Verbatim}[fontsize=\footnotesize]
  On branch master
  nothing to commit, working directory clean
\end{Verbatim}


\newpage
\textbf{Once you have completed all of the steps on the previous page, copy
the output of:\\
~\\
\texttt{$>$ git log -n2 -p | head -n 40}}

\vspace{45em}
\subsection{Preparing for the future}
\noindent
This question summarizes the setup and simple usage loop for git. Make some
changes, add them, and commit them. If you need to delete a file, use
\texttt{git~rm~file\_to\_delete}. From this moment onward, you should use git
for all of your EECS projects. Seriously. Use it.

Set up git repositories for all of your other EECS projects this term. If you
are currently in the middle of a project, create an initial commit at its
exact state right now (even though it is not yet finished / not working).
\textbf{\large \color{red}
  \begin{center}\texttt{* * *}\end{center}
  We will do more with version control again after spring break. It is
  important that you have at least one project that you have been using git
  with by then.
  \begin{center}\texttt{* * *}\end{center}
}
\noindent
Repositories from group projects are fine. If you use a different version
control system (\texttt{svn, bzr, hg, darcs, cvs, rcs, \dots}), choose one
project to manage with git instead, it's worth learning.

\newpage
\section{Pulling some pieces together}

Try the following:

\begin{lstlisting}
# The /tmp directory holds temporary files. It's a great place to do some
# quick tests or experiments, however it is automatically emptied every time
# the machine reboots. This means if your machine crashes, you *will* lose
# everything in /tmp.
#
# Don't ever put anything you care about in /tmp !!

> mkdir /tmp/throwaway
> cd /tmp/throwaway
> git init
> echo hello > world
> git add world
> git commit
\end{lstlisting}

\noindent
Git automatically opens a text editor for you to enter a commit message.

\bigskip
\noindent
\textbf{What text editor did git open for you? \rule{5cm}{0.1mm}}

\bigskip
\noindent
Git is not the only program that will open an editor. If you don't like the
default choice, we can change it using the environment variable
\texttt{EDITOR}. Try the following:

\begin{lstlisting}
> echo "it's a small" > world
> git add world
> export EDITOR="emacs -nw"
> git commit
\end{lstlisting}

\bigskip
\noindent
\textbf{What text editor did git open for you? \rule{5cm}{0.1mm}}

\bigskip
\noindent
\textbf{What does the \texttt{-nw} flag do?
  {\footnotesize (Hint: Try it without \texttt{-nw})}
}
\vspace{2cm}


\noindent
Perhaps, however, you do not want to change the editor used for every program,
but only for git. We can do that to. Try the following:

\begin{lstlisting}
> echo "on top of the" > world
> GIT_EDITOR=gedit git commit world
\end{lstlisting}

\noindent
\textbf{If you set \texttt{EDITOR} and \texttt{GIT\_EDITOR} to different
values, which takes priority?}
\vspace{2cm}

\noindent
\textbf{Describe how to ``permanently'' change the default \texttt{EDITOR} for
  your system to emacs
  {\footnotesize (Hint: What files did Q1 from week 2 read on startup?)}
}
\vspace{2cm}



\end{document}
